\section{Method}
\label{sec:method}

% Evidence: Models.md; models.py (BasicBlock / PreActBasicBlock / CifarResNet / CifarPreActResNet);
% experiments/configs/phase3|phase5; research_os/paper/canonical/data.json; G4-approved scope.
% Missing at write time: research_os/paper/outline.md, claim_ledger.md, citation_plan.md.

We study depth-20 CIFAR residual networks in the $6n{+}2$ family with $n{=}3$
(three stages with $k{=}(20{-}2)/6{=}3$ residual units per stage).
All residual models below use \emph{shortcut Option~A} at spatial/channel transitions:
stride-2 max-pooling on the identity path and zero-padding of new channels
(no learned $1{\times}1$ projection).
Plain networks omit the shortcut.
Unless noted, convolutional filters are $3{\times}3$.

\subsection{Baselines: Plain-20 and ResNet-20-A}
\label{sec:method-baselines}

\paragraph{Plain-20.}
A non-residual convolutional network with the same depth and stage widths as the residual counterpart,
but without identity shortcuts.
In our implementation, each basic unit follows a post-activation layout on the residual branch
without an additive skip:
$\mathrm{Conv}{\rightarrow}\mathrm{BN}{\rightarrow}\mathrm{ReLU}{\rightarrow}\mathrm{Conv}{\rightarrow}\mathrm{BN}{\rightarrow}\mathrm{Add}{\rightarrow}\mathrm{ReLU}$
with the shortcut disabled (equivalently, no residual addition).
This model is the non-residual reference used to measure the benefit of residual learning
at depth~20.

\paragraph{ResNet-20-A (frozen control).}
The primary baseline replaces the plain unit with a post-activation residual block
\[
\mathrm{Conv}{\rightarrow}\mathrm{BN}{\rightarrow}\mathrm{ReLU}{\rightarrow}\mathrm{Conv}{\rightarrow}\mathrm{BN}{\rightarrow}\mathrm{Add}{\rightarrow}\mathrm{ReLU},
\]
where $\mathrm{Add}$ merges the residual branch with an Option~A identity path, followed by ReLU.
ResNet-20-A is \emph{frozen} for all subsequent improvement and ablation runs:
later experiments reuse its training recipe and do not retrain this checkpoint as a moving target.
Under the protocol in Section~\ref{sec:setup}, its best test error is $\ErrResNet\%$
(canonical snapshot; SoT manifests/results).

\subsection{Improvement: PreAct-ResNet-20-A}
\label{sec:method-preact}

The improvement question is whether replacing the post-activation residual unit by a
pre-activation unit---while keeping dataset, optimizer, schedule, Option~A, and depth fixed---reduces
test error relative to ResNet-20-A.
The full pre-activation residual branch is
\[
\mathrm{BN}{\rightarrow}\mathrm{ReLU}{\rightarrow}\mathrm{Conv}{\rightarrow}\mathrm{BN}{\rightarrow}\mathrm{ReLU}{\rightarrow}\mathrm{Conv},
\]
followed by $\mathrm{Add}$ with the Option~A identity path and \emph{no} post-addition ReLU
(ReLU is applied at the start of the residual transform / next block entry).
Stem and classifier head are kept compatible with the baseline network; the intentional change is
the residual-block activation order.
We denote this model \textbf{PreAct-ResNet-20-A}.
Its best test error under the frozen recipe is $\ErrPreAct\%$
(canonical snapshot).

\subsection{Ablations of the pre-activation unit}
\label{sec:method-ablations}

To localize the contribution of the two $\mathrm{BN}{+}\mathrm{ReLU}$ stages inside the full PreAct unit,
we define two controlled ablations that share Option~A and the same frozen training recipe:

\paragraph{Ablation~A.}
Remove the \emph{first} $\mathrm{BN}{+}\mathrm{ReLU}$.
The residual branch becomes
$\mathrm{Conv}{\rightarrow}\mathrm{BN}{\rightarrow}\mathrm{ReLU}{\rightarrow}\mathrm{Conv}{\rightarrow}\mathrm{Add}$.
Best test error: $\ErrAblationA\%$ (canonical snapshot).

\paragraph{Ablation~B.}
Remove the \emph{second} $\mathrm{BN}{+}\mathrm{ReLU}$.
The residual branch becomes
$\mathrm{BN}{\rightarrow}\mathrm{ReLU}{\rightarrow}\mathrm{Conv}{\rightarrow}\mathrm{Conv}{\rightarrow}\mathrm{Add}$.
In our runs this configuration did not yield a comparable accuracy trajectory:
training failed with numerical instability (\AblationBStatus; no best-test-error baseline is reported).

\subsection{What is held fixed vs.\ what varies}
\label{sec:method-controls}

Across PreAct and both ablations, configs prescribe the same CIFAR-10 split, depth-20 width pattern,
Option~A rule, and the post-activation training recipe used for ResNet-20-A
(Section~\ref{sec:setup}).
The sole intentional architectural factor under study is the residual-block nonlinearity/BN placement
(full PreAct vs.\ dropping the first or second pre-activation stage).
We do not introduce attention, Squeeze-and-Excitation, Cutout, Mixup, dropout, or other
regularizers in these runs.
We do \emph{not} claim an exhaustive log-level audit that every Phase~3/5 trajectory matched the
shared recipe with zero silent augmentation drift
\texttt{[NEEDS\_VERIFICATION]}.
